% TOP_PROBLEM: {"id": 30006789, "title": "Kaplansky's Zero-Divisor Conjecture for Torsion-Free Groups", "category_id": 4, "category": {"id": 4, "name": "algebra", "display_name": "Algebra", "description": "Group theory, ring theory, field theory, and algebraic structures.", "slug": "algebra", "order_index": 4, "created_at": "2026-07-31T15:26:25.671Z"}, "status": "open"}
% FIRST_POSTED: 2026-10-06
% ATTEMPT_STATUS: solved
% OPENAI_CLAIM: {"schema_version":1,"source_repo":"https://github.com/openai/math","source_commit":"adc7f1241b42e322a6451854ab7e4b4c146bf78a","release_date":"2026-10-06","match":"full","resolution":"disproved","independently_reviewed":false,"second_pass":"agreed","families":[{"family":"196","title":"A counterexample to Kaplansky’s zero-divisor conjecture","subject":"Algebra","manuscripts":[{"title":"A Torsion-Free Group Algebra with Zero Divisors","pdf_path":"preprints/A-Torsion-Free-Group-Algebra-with-Zero-Divisors-September-23-2026/paper.pdf","readme_path":"preprints/A-Torsion-Free-Group-Algebra-with-Zero-Divisors-September-23-2026/README.md","date":"2026-09-23","theorem_ref":"Theorem 1.1"}],"lean":{"doc_path":"lean/docs/196.md","comparators":["lean/ComparatorChallenges/TorsionFreeZeroDivisors.json","lean/ComparatorChallenges/TorsionFreeZeroDivisors.lean"],"declarations":["OAI.TorsionFreeZeroDivisors.main"],"covers_main_theorem":true},"reasoning_trace":null},{"family":"197","title":"A torsion-free group algebra that is not directly finite","subject":"Algebra","manuscripts":[{"title":"A Torsion-Free Group Algebra That Is Not Directly Finite","pdf_path":"preprints/A-Torsion-Free-Group-Algebra-That-Is-Not-Directly-Finite-October-4-2026/direct-finiteness.pdf","readme_path":"preprints/A-Torsion-Free-Group-Algebra-That-Is-Not-Directly-Finite-October-4-2026/README.md","date":"2026-10-04","theorem_ref":"Theorem 1.1"}],"lean":{"doc_path":"lean/docs/197.md","comparators":[],"declarations":[],"covers_main_theorem":false},"reasoning_trace":"reasoning_traces/kaplansky-direct-finiteness-characteristic-two.pdf"}]}
\documentclass[11pt]{article}
\usepackage[margin=25mm]{geometry}
\usepackage{fontspec}
\setmainfont{Latin Modern Roman}
\usepackage{amsmath,amssymb,mathtools}
\usepackage{unicode-math}
\setmathfont{Latin Modern Math}
\usepackage{enumitem}
\usepackage{xurl,hyperref,bookmark}
\hypersetup{colorlinks=true,urlcolor=blue}
\setcounter{secnumdepth}{0}
\setlength{\parindent}{0pt}
\setlength{\parskip}{5pt}
\setlength{\emergencystretch}{3em}
\begin{document}
\section*{Kaplansky's Zero-Divisor Conjecture for Torsion-Free Groups}

\subsection{OpenAI's claim}
The linked manuscripts present the following claims by OpenAI: A finitely presented torsion-free group is claimed to have nonzero elements $a,c\in\mathbb F_2[G]$ with $ac=0$, despite admitting a finite two-dimensional classifying complex; Another construction claims a torsion-free, finitely presented $G$ for which $\mathbb F_2[G]$ contains $a,b$ satisfying $ab=1$ and $ba\ne1$. Together they give a claimed negative answer in the matching scope.

\subsection{Outline of the argument}
Finite labeled graphs encode supports whose product coefficients cancel modulo two. A random-label counting argument excludes planar path arrangements; cone surgeries turn that exclusion into asphericity and root separation, so the coefficients remain nonzero in a torsion-free group.

An unequal-side parity pattern supplies the one-sided inverse relation in labeled graphs. Weighted counts handle repeated edge traversals, and a separator reduces arbitrarily large planar arrangements to bounded path collections; cone attachment and root separation then preserve the inverse defect while proving torsion-freeness.

\subsection{Scope relative to this problem}
Theorem 1.1 gives torsion-free G and nonzero \(\alpha\),\(\beta\)\(\in\)F\({}_{2}\)[G] with \(\alpha\)\(\beta\)=0, so a single allowed field disproves the record’s all-fields assertion. Theorem 1.1 supplies torsion-free G with ac=0 and c\(\neq\)0, while ab=1 forces a\(\neq\)0, giving the scalar zero-divisor counterexample required to negate the universal assertion. Within these assumptions, this is a claimed negative answer to the stated question.

\subsection{Formal verification}
For family 196, the documented Lean theorem includes the torsion-free characteristic-two zero product, finite presentation, and a finite two-dimensional classifying model. The listed Comparator files are \texttt{TorsionFreeZeroDivisors.lean}. Named declarations include \texttt{OAI.TorsionFreeZeroDivisors.main}. For family 197, the documented Lean entries include the scalar inverse defects in the specified characteristics and an operator-invertible determinant counterexample. The characteristic-two search recipe is included, but the nonsofic consequence, the extra spectral-integral conclusion, and the later torsion-free construction are outside the selected statements. No Comparator is assigned to the selected manuscript here; the family page selects other papers or supporting statements. This site has not executed the Lean code. The scope descriptions do not certify the remaining manuscript arguments.

\subsection{Status caveat}
OpenAI attributes these manuscripts to work by an internal model. Its release cautions that arguments without formalization can contain mistakes. This catalogue has not independently checked the proofs; the labels report the mathematical claims attributed to OpenAI.

\subsection{Sources}
\begin{itemize}
\item OpenAI, \emph{A Torsion-Free Group Algebra with Zero Divisors}, 2026-09-23; citation key \texttt{OAI:A-Torsion-Free-Group-Algebra-with-Zero-Divisors-September-23-2026}. \href{https://github.com/openai/math/blob/adc7f1241b42e322a6451854ab7e4b4c146bf78a/preprints/A-Torsion-Free-Group-Algebra-with-Zero-Divisors-September-23-2026/paper.pdf}{PDF (GitHub)}; \href{https://github.com/openai/math/raw/adc7f1241b42e322a6451854ab7e4b4c146bf78a/preprints/A-Torsion-Free-Group-Algebra-with-Zero-Divisors-September-23-2026/paper.pdf}{PDF (direct)}; \href{https://github.com/openai/math/blob/adc7f1241b42e322a6451854ab7e4b4c146bf78a/preprints/A-Torsion-Free-Group-Algebra-with-Zero-Divisors-September-23-2026/README.md}{Manuscript page and citation}; \href{https://github.com/openai/math/blob/main/preprints/A-Torsion-Free-Group-Algebra-with-Zero-Divisors-September-23-2026/paper.pdf}{Latest version}.
\item Family 196: \href{https://github.com/openai/math/blob/adc7f1241b42e322a6451854ab7e4b4c146bf78a/lean/docs/196.md}{Lean scope}.
\item \href{https://github.com/openai/math/blob/adc7f1241b42e322a6451854ab7e4b4c146bf78a/lean/ComparatorChallenges/TorsionFreeZeroDivisors.json}{Comparator TorsionFreeZeroDivisors.json}.
\item \href{https://github.com/openai/math/blob/adc7f1241b42e322a6451854ab7e4b4c146bf78a/lean/ComparatorChallenges/TorsionFreeZeroDivisors.lean}{Comparator TorsionFreeZeroDivisors.lean}.
\item OpenAI, \emph{A Torsion-Free Group Algebra That Is Not Directly Finite}, 2026-10-04; citation key \texttt{OAI:A-Torsion-Free-Group-Algebra-That-Is-Not-Directly-Finite-October-4-2026}. \href{https://github.com/openai/math/blob/adc7f1241b42e322a6451854ab7e4b4c146bf78a/preprints/A-Torsion-Free-Group-Algebra-That-Is-Not-Directly-Finite-October-4-2026/direct-finiteness.pdf}{PDF (GitHub)}; \href{https://github.com/openai/math/raw/adc7f1241b42e322a6451854ab7e4b4c146bf78a/preprints/A-Torsion-Free-Group-Algebra-That-Is-Not-Directly-Finite-October-4-2026/direct-finiteness.pdf}{PDF (direct)}; \href{https://github.com/openai/math/blob/adc7f1241b42e322a6451854ab7e4b4c146bf78a/preprints/A-Torsion-Free-Group-Algebra-That-Is-Not-Directly-Finite-October-4-2026/README.md}{Manuscript page and citation}; \href{https://github.com/openai/math/blob/main/preprints/A-Torsion-Free-Group-Algebra-That-Is-Not-Directly-Finite-October-4-2026/direct-finiteness.pdf}{Latest version}.
\item Family 197: \href{https://github.com/openai/math/blob/adc7f1241b42e322a6451854ab7e4b4c146bf78a/lean/docs/197.md}{Lean scope}.
\item Family 197: \href{https://github.com/openai/math/blob/adc7f1241b42e322a6451854ab7e4b4c146bf78a/reasoning\_traces/kaplansky-direct-finiteness-characteristic-two.pdf}{Reasoning summary}.
\item \href{https://github.com/openai/math/blob/adc7f1241b42e322a6451854ab7e4b4c146bf78a/CONTENTS.md}{OpenAI catalogue entry}.
\end{itemize}
\end{document}
